% GENERATED FILE. Edit canonical lessons, then run npm run generate:books.
% canonical-source-hash: fnv1a64:5473039a85e24d66
% canonical-lessons: LA-C378-exerceo, LA-C378-fabricor, LA-C378-texo, LA-C378-suo, LA-C378-sculpo

\chapter{Making Things by Hand}
\label{ch:la-making-things-by-hand}
\hlchaptermodality{378}

\begin{chapteropening}
\textbf{By the end of this chapter:} \emph{I can say that someone practises a craft, and that he makes, weaves, sews or carves something.}

Complete the last lesson of chapter 378: I can say that someone practises a craft, and that he makes, weaves, sews or carves something.
\end{chapteropening}

\section[exerceō, exercēre]{exerceō, exercēre — to train, to practise; to keep busy}

\label{lesson:LA-C378-exerceo}



\begin{quote}
Before the new one: say the Latin for a mason, then the Latin for a sculptor of statues.
\end{quote}

\subsection*{You'll want to know: exerceō, exercēre}
\textbf{exerceō, exercēre} — \textquotedblleft{}to train, to practise; to keep busy\textquotedblright{}.

\begin{quote}
Pistor artem suam multōs annōs exercuit.
\end{quote}

\begin{quote}
\emph{The baker practised his trade for many years.}
\end{quote}

\emph{exerceō, exercēre, exercuī, exercitum}. \emph{Artem exercēre}, to practise a craft; \emph{sē exercēre}, to train. English \emph{exercise}.

\begin{cousinweb}
\emph{exerceō} is built from \emph{ex} and \emph{arceō} \textquotedblleft{}to shut in, to keep off\textquotedblright{} (EtymDB; Lewis \& Short): to drive animals out to work, then to keep anyone busy. English \emph{exercise} comes from it.
\end{cousinweb}

\begin{grammarlens}[title={What you've built}]
One more word for telling, explaining and describing at B1.
\end{grammarlens}

\subsection*{Your turn}
\begin{itemize}
\raggedright
  \item \emph{Say it:} \emph{exerceō, exercēre}
  \item \emph{Say it:} \emph{exerceō, exercēre}, once more
  \item \emph{Say it:} a sentence of your own with \emph{exerceō, exercēre}
  \item \emph{Recall:} say the Latin for a mason, then the Latin for a sculptor of statues, then say \emph{exerceō, exercēre} again
\end{itemize}

\begin{tcolorbox}[breakable,colback=teal!4,colframe=teal!35!black,title={Before you move on}]
What does exerceō, exercēre mean? (\textquotedblleft{}To train, to practise; to keep busy\textquotedblright{}.) Say it once more.
\end{tcolorbox}
\section[fabricor, fabricārī]{fabricor, fabricārī — to make, to build, to forge}

\label{lesson:LA-C378-fabricor}



\begin{quote}
Before the new one: say the Latin for a sculptor of statues, then the Latin for to train.
\end{quote}

\subsection*{You'll want to know: fabricor, fabricārī}
\textbf{fabricor, fabricārī} — \textquotedblleft{}to make, to build, to forge\textquotedblright{}.

\begin{quote}
Faber gladium fabricātur.
\end{quote}

\begin{quote}
\emph{The smith is forging a sword.}
\end{quote}

Deponent: \emph{fabricor, fabricārī, fabricātus sum}, built on \emph{fabrica}, a workshop, and \emph{faber}, a smith. English \emph{fabricate}.

\begin{cousinweb}
\emph{fabricor} is built on \emph{fabrica} \textquotedblleft{}a workshop, a craft\textquotedblright{} (Lewis \& Short), which comes from \emph{faber} \textquotedblleft{}a smith, a craftsman\textquotedblright{}. English \emph{fabricate} comes from it.
\end{cousinweb}

\begin{grammarlens}[title={What you've built}]
One more word for telling, explaining and describing at B1.
\end{grammarlens}

\subsection*{Your turn}
\begin{itemize}
\raggedright
  \item \emph{Say it:} \emph{fabricor, fabricārī}
  \item \emph{Say it:} \emph{fabricor, fabricārī}, once more
  \item \emph{Say it:} a sentence of your own with \emph{fabricor, fabricārī}
  \item \emph{Recall:} say the Latin for a sculptor of statues, then the Latin for to train, then say \emph{fabricor, fabricārī} again
\end{itemize}

\begin{tcolorbox}[breakable,colback=teal!4,colframe=teal!35!black,title={Before you move on}]
What does fabricor, fabricārī mean? (\textquotedblleft{}To make, to build, to forge\textquotedblright{}.) Say it once more.
\end{tcolorbox}
\section[texō, texere]{texō, texere — to weave; to build}

\label{lesson:LA-C378-texo}



\begin{quote}
Before the new one: say the Latin for to train, then the Latin for to make.
\end{quote}

\subsection*{You'll want to know: texō, texere}
\textbf{texō, texere} — \textquotedblleft{}to weave; to build\textquotedblright{}.

\begin{quote}
Fēmina vestem ē lānā texit.
\end{quote}

\begin{quote}
\emph{The woman weaves a garment from wool.}
\end{quote}

\emph{texō, texere, texuī, textum}. A speech too can be woven. English \emph{text}, \emph{textile} and \emph{texture}.

\begin{cousinweb}
\emph{texō} is inherited from Proto-Italic \emph{*teksō} (EtymDB), from a root meaning \textquotedblleft{}to fashion, to build\textquotedblright{} that Greek \emph{tektōn} \textquotedblleft{}a builder\textquotedblright{} shares. English \emph{text} and \emph{textile} come from it.
\end{cousinweb}

\begin{grammarlens}[title={What you've built}]
One more word for telling, explaining and describing at B1.
\end{grammarlens}

\subsection*{Your turn}
\begin{itemize}
\raggedright
  \item \emph{Say it:} \emph{texō, texere}
  \item \emph{Say it:} \emph{texō, texere}, once more
  \item \emph{Say it:} a sentence of your own with \emph{texō, texere}
  \item \emph{Recall:} say the Latin for to train, then the Latin for to make, then say \emph{texō, texere} again
\end{itemize}

\begin{tcolorbox}[breakable,colback=teal!4,colframe=teal!35!black,title={Before you move on}]
What does texō, texere mean? (\textquotedblleft{}To weave; to build\textquotedblright{}.) Say it once more.
\end{tcolorbox}
\section[suō, suere]{suō, suere — to sew, to stitch}

\label{lesson:LA-C378-suo}



\begin{quote}
Before the new one: say the Latin for to make, then the Latin for to weave.
\end{quote}

\subsection*{You'll want to know: suō, suere}
\textbf{suō, suere} — \textquotedblleft{}to sew, to stitch\textquotedblright{}.

\begin{quote}
Sūtor calceōs suit.
\end{quote}

\begin{quote}
\emph{The shoemaker stitches shoes.}
\end{quote}

\emph{suō, suere, suī, sūtum}. \emph{Sūtor}, a shoemaker, is the man who sews. English \emph{suture} comes from it; English \emph{sew} is a native cousin, not a loan.

\begin{cousinweb}
\emph{suō} is inherited from Proto-Italic \emph{*sūō} (EtymDB), from the same Proto-Indo-European root as English \emph{sew}, which came down through Old English, not from Latin.
\end{cousinweb}

\begin{grammarlens}[title={What you've built}]
One more word for telling, explaining and describing at B1.
\end{grammarlens}

\subsection*{Your turn}
\begin{itemize}
\raggedright
  \item \emph{Say it:} \emph{suō, suere}
  \item \emph{Say it:} \emph{suō, suere}, once more
  \item \emph{Say it:} a sentence of your own with \emph{suō, suere}
  \item \emph{Recall:} say the Latin for to make, then the Latin for to weave, then say \emph{suō, suere} again
\end{itemize}

\begin{tcolorbox}[breakable,colback=teal!4,colframe=teal!35!black,title={Before you move on}]
What does suō, suere mean? (\textquotedblleft{}To sew, to stitch\textquotedblright{}.) Say it once more.
\end{tcolorbox}
\section[sculpō, sculpere]{sculpō, sculpere — to carve, to engrave}

\label{lesson:LA-C378-sculpo}



\begin{quote}
Before the new one: say the Latin for to weave, then the Latin for to sew.
\end{quote}

\subsection*{You'll want to know: sculpō, sculpere}
\textbf{sculpō, sculpere} — \textquotedblleft{}to carve, to engrave\textquotedblright{}.

\begin{quote}
Statuārius nōmen in saxō sculpsit.
\end{quote}

\begin{quote}
\emph{The sculptor carved the name in stone.}
\end{quote}

\emph{sculpō, sculpere, sculpsī, sculptum}. English \emph{sculpture} and \emph{sculptor}.

\begin{cousinweb}
\emph{sculpō} is a form beside \emph{scalpō} \textquotedblleft{}to scratch, to scrape\textquotedblright{} (EtymDB; Lewis \& Short). English \emph{sculpture} and \emph{sculptor} come from it.
\end{cousinweb}

\begin{grammarlens}[title={What you've built}]
One more word for telling, explaining and describing at B1.
\end{grammarlens}

\subsection*{Your turn}
\begin{itemize}
\raggedright
  \item \emph{Say it:} \emph{sculpō, sculpere}
  \item \emph{Say it:} \emph{sculpō, sculpere}, once more
  \item \emph{Say it:} a sentence of your own with \emph{sculpō, sculpere}
  \item \emph{Recall:} say the Latin for to weave, then the Latin for to sew, then say \emph{sculpō, sculpere} again
\end{itemize}

\begin{tcolorbox}[breakable,colback=teal!4,colframe=teal!35!black,title={Before you move on}]
What does sculpō, sculpere mean? (\textquotedblleft{}To carve, to engrave\textquotedblright{}.) Say it once more.
\end{tcolorbox}
